This project investigated automatic classification of GitHub issues into bug, feature, and question categories using the NLBSE'24 dataset. The implementation compared classical sparse classifiers with pretrained Transformer and Sentence Transformer approaches within a shared data and evaluation pipeline.

The experiments show that the issue body contributes substantial information beyond the title alone. The TF-IDF + Logistic Regression baseline reaches 0.7636 Macro-F1 under repository-specific CV but decreases to 0.5876 under LOO, illustrating the difficulty of transfer to a repository that is absent from training. Under repository-specific CV, SetFit MPNet reaches 0.7953 and the five-iteration SetFit configuration reaches 0.7976. RoBERTa LoRA reaches 0.7994 under pooled CV; this result is reported separately because the pooled protocol uses a different training composition.

The title-only LOO extension reaches 0.5075. Its paired difference from title plus body supports a positive body-text contribution under the fixed-prediction bootstrap, with the largest descriptive gain on bugs. P5's illustrative LR--MPNet hard vote reaches 0.7885, below MPNet alone; combining models does not automatically improve performance. These are training-only, post-hoc analyses rather than official-test validation.

The project also highlights the importance of interpreting predictive performance together with evaluation protocol and computational cost. Sparse classifiers remain simple and inexpensive reference systems, while pretrained dense models provide stronger observed CV performance in the reported experiments. The SetFit iteration results further show that substantially longer training does not necessarily produce a better score.

The shared experiment pipeline provides a useful foundation for future extensions because it keeps data identity, model configuration, predictions, metrics, and available resource information together. Possible improvements include evaluating more pretrained models under LOO, repeating stochastic experiments with a consistent set of random seeds, collecting resource measurements on a single controlled hardware environment, and reporting additional per-class or error-analysis results. These extensions would strengthen the comparison while preserving the practical scope of the project report.
